%!TEX TS-program = xelatex
%!TEX encoding = UTF-8 Unicode

\documentclass[11pt, a4paper]{awesome-cv}

\geometry{left=1.4cm, top=.8cm, right=1.4cm, bottom=1.8cm, footskip=.5cm}

\colorlet{awesome}{awesome-red}

\setbool{acvSectionColorHighlight}{true}

\renewcommand{\acvHeaderSocialSep}{\quad\textbar\quad}


%-------------------------------------------------------------------------------
% PERSONAL INFORMATION
%-------------------------------------------------------------------------------
\name{}{김송아}
\position{Data Strategist \& Analyst}
\address{서울시 서초구 효령로 391}
\mobile{010-3493-8659}
\email{kimsonga17@gmail.com}
\github{kimsonga17}
\linkedin{kimsonga17}

\quote{``데이터 기반 의사결정 환경을 설계하고 조직 전반의 데이터 활용 체계를 구축해온 데이터 전략 및 분석 전문가입니다."}


%-------------------------------------------------------------------------------
\begin{document}

\makecvheader[C]

\makecvfooter
  {\today}
  {김송아~~~·~~~이력서}
  {\thepage}


%-------------------------------------------------------------------------------
% CV/RESUME CONTENT
%-------------------------------------------------------------------------------

\cvsection{Objectives}

\begin{itemize}
    \item 데이터 기반 의사결정이 가능한 환경을 설계하고, 조직 전반의 데이터 활용 체계를 구축해온 데이터 전략 및 분석 전문가입니다.
    \item 다양한 도메인의 데이터를 통합해 KPI를 정의하고, 유저 퍼널 및 행동 기반 분석을 통해 실질적인 서비스 개선과 전략 수립에 기여해왔습니다.
    \item 최근에는 블록체인 기반 신규 플랫폼에서 MVP 단계부터 전사 데이터 전략 수립, 로그 설계, 퍼널 분석 체계 구축까지 리딩하며 데이터 기반 문제 해결력을 강화했습니다.
    \item 개발, 기획, 사업, 마케팅, CS 등 다양한 조직과 협업해 유실 로그 보완, 도메인 간 데이터 연계, 공통 지표 수립 등을 이끌었습니다.
    \item 주간 데이터 그로스 미팅을 통해 실무 조직 간의 협업 구조를 정착시키고, 로그 개선과 인사이트 공유를 통해 데이터 리터러시와 실행력을 높이는 문화를 만들어 왔습니다.
\end{itemize}

\cvsection{Experience}

\begin{cventries}
  \cventry
    {Data Lead}
    {LINE Next}
    {경기 성남}
    {2024.04 -- 현재}
    {}

  \cventry
    {DOSI Strategy}
    {LINE Next}
    {경기 성남}
    {2022.10 -- 2024.04}
    {
    \begin{itemize}
        \item 글로벌 B2C-C2C NFT 플랫폼(DOSI)의 사업 전략 수립 과정에서 데이터 기반 인사이트를 제공하며 핵심 역할 수행, LINE NEXT 데이터팀 구성 및 설립 주도
        \item C-Level과 협력하여 핵심 비즈니스 및 제품 전략 개발을 위한 데이터 전략 총괄 및 인사이트 지원
        \item 5개 주요 서비스의 도메인 데이터를 통합 수집할 수 있도록 데이터 정의 체계를 표준화
        \item 공통 데이터 웨어하우스로의 파이프라인 개설을 주도하고, 분석 가능한 구조로 통합
        \item 데이터 표준화를 통해 수집 프로세스 일관성 확보 및 분석 목적에 최적화된 파이프라인 설계 주도, 전사 공통 KPI 리포팅 체계 기반 마련
        \item 데이터 리터러시 향상 및 협업 체계 구축을 위해 ‘주간 데이터 그로스 미팅’을 주도, 개발, 기획, 테크서포트 등 유관 조직과 데이터 기반 이슈를 공유하고 협업하는 구조를 정례화
      \end{itemize}
    }

  \cventry
    {CPO Strategy}
    {LINE Plus Corp}
    {경기 성남}
    {2020.12 -- 2022.10}
    {
      \begin{itemize}
        \item JP/TH/TW의 주요 메신저 플랫폼 CPO 전속 전략팀 소속으로, MAU 1.8억 플랫폼과 기타 Family Service(18개 가량) 전략관련 데이터 전반 서포트
        \item 데이터 분석 환경 구축, BI 도입 및 마트 개발
        \item KPI report / Insight Report 등 시각화 및 대시보드 제작
        \item Cross Platform Analysis - User activity, User retention, Service Overlap, AUM 등 서비스의 주요 KPI 개발
        \item FINTECH 서비스(LINE app, Pay, Wallet, 증권, 대출, 가상자산 거래소, 인터넷은행 등) KPI, PL 모니터링
        \item 핀테크 및 Blockchain 관련 업계동향 리서치
        \item LINE App Service 관련 Annual 회고 리포트 작성
      \end{itemize}
    }

  \cventry
    {User Management \& Data Analyst, BITFRONT VME Service Team}
    {LINE Biz Plus, LINE Financial Plus}
    {경기 판교}
    {2018.06 -- 2020.11}
    {
      \begin{itemize}
        \item RFM 및 User Segmentation 분석으로 회원 등급제 고안, 개인화 서비스 제안
        \item Promotion Scheme 분석 및 기획: 데이터 분석을 통해서 프로모션 구조 기획. New user 획득에 기여(MAU 대비 10\%)
        \item Loyalty Program 기획: 경쟁사 동향 리서치를 통해, 서비스 적용 제안
        \item User Lifecycle management 기획, 가입부터 이탈까지 유저행동에 대해 구조화
        \item User retention rate을 높이기 위한 Performance Marketing Project 리드
        \item 컨텐츠 기획을 통해 User 유입 지표 개선(10\%)
        \item 서비스 오픈 멤버로 Customer Support 전반 기획 및 프로세스 셋업
        \item KYC(Know Your Customer) 프로세스 셋업 및 개선
        \item CRM 툴 (Salesforce) 서비스 연동 및 프로세스 개선 작업 참여
        \item 고객센터 KPI, SLA 등 주요지표 수립
      \end{itemize}
    }

  \cventry
    {CRM Manager \& Data Analyst, Product Team}
    {KORBIT}
    {서울}
    {2016.09 -- 2018.06}
    {
      \begin{itemize}
        \item 데이터 기반 VOC 리포트 고객의 주요 피드백 선별, 유저 실적 정보와 결합하여 위클리 리포트 기획
        \item RFM 분석, 이탈회원 모델링 등 고객가치측정 분석
        \item 회사 매출 및 거래 자료 시각화, 전사 공유
        \item 검색어 광고 퍼포먼스 관리 및 측정
        \item 마케팅 효과 분석 리포트 작성
        \item CRM툴 도입: CS team 업무 효율성 높임으로써 고객만족도 신장
        \item 아웃소싱 컨택 센터 오픈 프로젝트 리드
        \item Customer Support 팀 KPI 및 효율지표 개발
        \item 고객 문의사항 분류 및 빈도 분석: 대시보드 생성, 고객 빈출 문의사항 도출하여 서비스 개선 방향 제안
      \end{itemize}
    }

  \cventry
    {Researcher}
    {Hanyang University}
    {서울}
    {2015.08 -- 2016.03}
    {
      \begin{itemize}
        \item 유학생 심리상담 및 분석
        \item 유학생 심리검사 및 분석
        \item 유학생 생활 적응 조사 및 분석
      \end{itemize}
    }

%   \cventry
%     {Researcher}
%     {Ehwa Diamond}
%     {경기 오산}
%     {2011.08 -- 2012.07}
%     {
%       \begin{itemize}
%         \item 바이어 관리(영국)
%         \item 영업 실적 관리
%         \item 의전
%         \item 재고관리
%       \end{itemize}
%     }
\end{cventries}

\cvsection{Skills}

\begin{itemize}
    \item R
    \item SQL
    \item Python
    \item Tableau
    \item GA4
    \item JIRA/CONFLUENCE
\end{itemize}

\cvsection{Education}

\begin{cventries}
  \cventry
    {석사, 심리학 상담전공}
    {가톨릭대학교}
    {경기}
    {2013년 9월 -- 2015년 8월}
    {
      \begin{itemize}
        \item 논문: 웅대성-취약성 자기애와 대인소망의 관계, 성인애착의 매개효과
      \end{itemize}
    }

  \cventry
    {학사, 문화인류학 전공, 국어국문학 부전공}
    {한양대학교}
    {경기}
    {2005년 3월 -- 2011년 2월}
    {}
\end{cventries}

%-------------------------------------------------------------------------------
\end{document}